%------------------------------------------------------------------------------%
\section{Limites e Continuidade}
\label{sec:limites_Continuidade}

A seguir apresentamos a definição formal do limite de uma função real.

%------------------------------------------------------------------------------%
\begin{definicao}{Limite de Função Real}{def:limite-funcao-real}
  Seja $f$ uma função definida em intervalo aberto que contém do ponto $a\in\R$,
  podendo não ser definida em $a$. Dizemos que o \emph{Limite}\index{Limite!função real}
  de $f(x)$ quando $x$ se aproxima de $a$ é $L\in\R$ e escrevemos
  \[
  \lim_{x\to a} f(x) = L
  \]
  se para cada $\varepsilon>0$ existir $\delta>0$ tal que
  \[
  0 < \abs{x-a} < \delta \qquad\Rightarrow\qquad \abs{f(x)-L} < \varepsilon
  \]
\end{definicao}
%------------------------------------------------------------------------------%

Essa definição pode ser estendida para o caso em que $a$ é $\pm\infty$.
Quando o limite existe ele possui as propriedades descritas na proposição a seguir.

%------------------------------------------------------------------------------%
\begin{proposicao}{Propriedades do Limite de Funções Reais}{pro:limite-funcao-real}
  Supondo que $c$ é uma constante, $f$ e $g$ são funções reais e que os limites
  \[
  \lim_{x\to a} f(x) \qquad\text{ e }\qquad \lim_{x\to a} g(x)
  \]
  existem, então valem as \emph{Propriedades do Limite}
  \begin{enumerate}[parsep=11pt]
    \item \( \lim_{x\to a} \Big[ f(x) + g(x) \Big] = \lim_{x\to a} f(x) + \lim_{x\to a} g(x) \)
    \item \( \lim_{x\to a} \Big[ f(x) - g(x) \Big] = \lim_{x\to a} f(x) - \lim_{x\to a} g(x) \)
    \item \( \lim_{x\to a} cf(x) = c \lim_{x\to a} f(x) \)
    \item \( \lim_{x\to a} \Big[ f(x)g(x) \Big] = \lim_{x\to a} f(x) \cdot \lim_{x\to a} g(x) \)
    \item \( \lim_{x\to a} \frac{f(x)}{g(x)} = \dfrac{\displaystyle\lim_{x\to a} f(x)}
    {\displaystyle\lim_{x\to a} d(x)} \)
    \qquad desde que \(\lim_{x\to a} d(x) \neq 0\)
  \end{enumerate}
\end{proposicao}
%------------------------------------------------------------------------------%

Uma propriedade de algumas funções reais é não possuir buracos ou saltos em seu
gráfico. Dizemos que essa funções são contínuas e a próxima definição nos fornece
uma caracterização para essas funções.

%------------------------------------------------------------------------------%
\begin{definicao}{Continuidade de Função Real}{def:continuidade-funcao-real}
  Uma função real é \emph{Contínua}\index{Função!contínua} em um ponto $a$ se
  \begin{enumerate}[parsep=9pt]
    \item \( f(a) \) existe
    \item \( \lim_{x\to a} f(x) \) existe
    \item \( \lim_{x\to a} f(x) = f(a) \)
  \end{enumerate}
\end{definicao}
%------------------------------------------------------------------------------%

Um resultado importante sobre funções contínuas é o Teorema do Valor Intermediário.

%------------------------------------------------------------------------------%
\begin{teorema}{Teorema do Valor Intermediário}{teo:valor-intermediario}
  Suponha\index{Teorema!valor intermediário}
  que a função $f$ é contínua em todos os pontos de um intervalo fechado
  $[a, b]$ e que \(f(a)\neq f(b)\). Assuma que $p$ é um valor entre $f(a)$ e $f(b)$
  então existe um número $c\in(a, b)$ tal que $f(c)=p$.
\end{teorema}
%------------------------------------------------------------------------------%


%------------------------------------------------------------------------------%
